% ---------------------------------------------------------------
%  Section V -- Methods
% ---------------------------------------------------------------

\subsection{Ingestion and Validation}

All external inputs are validated against Pydantic~v2 schemas before
downstream processing.
Four schema families are defined and enforced at the ingestion boundary:
\begin{itemize}
  \item \textit{SensorEvent}: device identifier, ISO-8601 timestamp,
        event type enumeration, payload dictionary (typed), and optional
        WGS-84 coordinates.
  \item \textit{GNSSPoint}: device identifier, POSIX timestamp, WGS-84
        latitude and longitude, speed (m\,s$^{-1}$), heading (degrees),
        and horizontal accuracy (m).
  \item \textit{WeatherSnapshot}: region identifier, acquisition timestamp,
        temperature (°C), precipitation (mm\,h$^{-1}$), wind speed
        (m\,s$^{-1}$), relative humidity (\%), and data status.
  \item \textit{PublicEvent}: event identifier, venue coordinates,
        ISO-8601 start and end timestamps, and estimated attendance.
\end{itemize}
Validation failures are logged with the offending field, the expected type,
and the received value; invalid records are discarded without halting the
pipeline.
Provenance records $\pi_s$ (Eq.~\ref{eq:provenance}) are attached immediately
after successful validation and before the merge step.
File-based ingestion from configured directories supports validated real
snapshots for reproducible offline evaluation; synthetic directories are
not used in any application-level code path.

\subsection{Digital Twin State Assembly}

Algorithm~\ref{alg:assembly} formalizes the assembler.
The assembler follows a merge-only policy: it attempts each configured source
in priority order and attaches the result (populated payload or empty payload
with an explicit status tag) to the snapshot provenance list $\Pi$.
No synthetic data is injected at any step, and no source output is
silently suppressed.

\begin{algorithm}[t]
\caption{Digital Twin State Assembly}
\label{alg:assembly}
\begin{algorithmic}[1]
\Require Configured source list $\mathcal{S}$, request timestamp $t$
\Ensure Snapshot $(G, t, \DS_{\mathrm{global}}, \Pi)$

\State $G \leftarrow (V{=}\emptyset,\; E{=}\emptyset,\; W{=}\emptyset)$;\;
       $\Pi \leftarrow [\;]$
\For{each source $s \in \mathcal{S}$ in priority order}
  \State $(\mathrm{data}_s,\; \DS_s) \leftarrow \textsc{Fetch}(s, t)$
  \Comment{Returns empty + status on timeout}
  \If{$\DS_s \in \{\texttt{live},\, \texttt{recorded}\}$}
    \State $G \leftarrow \textsc{Merge}(G,\; \mathrm{data}_s)$
    \Comment{Merge graph and dynamic state}
  \EndIf
  \State $\Pi \mathrel{+}= (s.\mathrm{name},\;\DS_s,\;t,\;\mathrm{checksum}(s))$
\EndFor
\State $\DS_{\mathrm{global}} \leftarrow \textsc{Aggregate}(\{\DS_s : s \in \mathcal{S}\})$
       \Comment{Rule: Eq.~(\ref{eq:global_status})}
\State \Return $(G,\; t,\; \DS_{\mathrm{global}},\; \Pi)$
\end{algorithmic}
\end{algorithm}

\subsection{Forecasting Baseline}

The forecasting module consumes the assembled snapshot $(G, t, \DS, \Pi)$.
In the heuristic baseline (active in the current release), the predicted
speed on edge $e$ at look-ahead $h$ time steps is
\begin{equation}
  \hat{s}_e(t+h) =
  \begin{cases}
    s_e(t) \cdot g(h) & \text{if } \DS_{\mathrm{traffic}} = \texttt{live}, \\
    s_e^0             & \text{otherwise,}
  \end{cases}
  \label{eq:heuristic_forecast}
\end{equation}
where $g(h) = e^{-\mu h}$ is an exponential decay towards free-flow speed
with decay rate $\mu > 0$ and $s_e^0 = \ell_e / \tau_e^0$ is the free-flow
speed.
The response inherits $\DS$ from the twin, ensuring that downstream consumers
can detect degraded input quality at any horizon.

The ST-GNN extension point follows Eqs.~(\ref{eq:forecast})--(\ref{eq:gcn}).
Training will proceed via mini-batch stochastic gradient descent on
$\mathcal{L}(\theta) = \mathrm{MAE}(\hat{\mathbf{Y}}, \mathbf{Y})$ with
the Adam optimizer (learning rate $\eta = 10^{-3}$, weight decay
$\lambda_w = 10^{-4}$, batch size $B = 64$) once a historical observation
store of sufficient length is assembled.
Table~\ref{tab:hyperparams} summarizes the target hyperparameter
configuration for the ST-GNN training phase.

\begin{table}[t]
\caption{Target hyperparameter configuration for the ST-GNN training phase.
         No training results are reported in the current release; these
         values define the planned experimental protocol.}
\label{tab:hyperparams}
\centering
\renewcommand{\arraystretch}{1.20}
\begin{tabular}{@{}lll@{}}
\toprule
\textbf{Hyperparameter} & \textbf{Symbol} & \textbf{Value} \\
\midrule
Lookback window        & $T$           & 12 steps (1 h at 5-min polling) \\
Forecast horizon       & $H$           & 12 steps (1 h ahead) \\
GNN hidden channels    & $d_h$         & 64 \\
GNN depth              & $L$           & 3 layers \\
Adjacency bandwidth    & $\sigma_\ell$ & \SI{500}{m} \\
GRU hidden size        & $d_\mathrm{gru}$ & 64 \\
Optimizer              & --            & Adam \\
Learning rate          & $\eta$        & $10^{-3}$ \\
Weight decay           & $\lambda_w$   & $10^{-4}$ \\
Batch size             & $B$           & 64 graphs \\
Maximum epochs         & --            & 100 \\
Early stopping patience& --            & 10 epochs \\
Train / val / test split & --          & 70\,/\,10\,/\,20\,\% \\
\bottomrule
\end{tabular}
\end{table}

\subsection{Routing}

The routing service solves Eq.~(\ref{eq:routing}) subject to
Eq.~(\ref{eq:routing_constraint}) over the assembled graph $G$.
When $G$ is empty or incomplete due to unconfigured sources, the service
returns a minimal response with the inherited data status and an explanatory
message in the \texttt{message} field.
The preference vector $[\alpha, \beta, \gamma, \delta]$ is accepted per
request; the default configuration minimizes travel time only
($\alpha=1$, $\beta=\gamma=\delta=0$).
OpenTripPlanner and Valhalla will realize full multi-modal routing once
GTFS feeds are available; the routing-engine abstraction layer
ensures that switching backends does not affect the upstream API contract
or the provenance protocol.

\subsection{Equity Computation}

The equity service reads district-level data from a path configured by
the environment variable \texttt{EQUITY\_DATA\_PATH}.
Each record in the \texttt{district\_scores.json} file provides indicator
values $\{x_{d,m}\}_{d,m}$.
The service computes $z_{d,m}$ via Eq.~(\ref{eq:zscore_district}),
the composite $\MES_d$ via Eq.~(\ref{eq:mes}), the equity gap
$\Delta_{\MES}$ via Eq.~(\ref{eq:equity_gap}), and the equity dispersion
$\mathrm{Var}(\MES)$ via Eq.~(\ref{eq:mes_var}).
The API response includes $\MES_d$ for all districts, $\Delta_{\MES}$,
$\mathrm{Var}(\MES)$, and a ranked list of the five least- and
most-accessible districts.
When \texttt{EQUITY\_DATA\_PATH} is not set, the service returns
$\DS = \texttt{unavailable}$ and empty payload fields.

\subsection{Anomaly Detection}

The anomaly pipeline evaluates Eq.~(\ref{eq:zscore_anomaly}) over edge-speed
signals using a configurable rolling window ($w = 12$ steps by default,
corresponding to one hour at five-minute polling intervals).
Severity is computed via Eq.~(\ref{eq:anomaly_severity}) when event and
weather signals are available; if not, the severity reduces to
$S_t = |z_t|$ (i.e., $\lambda_1 = 1$, $\lambda_2 = \lambda_3 = 0$).
The false-positive rate estimate (Eq.~\ref{eq:fpr}) is included in each
anomaly response to enable threshold calibration by operators.
Detected anomalies are labeled with: type (speed deviation, occupancy
spike, provider alert), severity
(\texttt{low}\,/\,\texttt{medium}\,/\,\texttt{high}\,/\,\texttt{critical}),
confidence (derived from Eq.~\ref{eq:fpr} at threshold $\tau_S$),
affected geography (edge set or bounding box), and source family
(\texttt{observed}, \texttt{inferred}, \texttt{provider-reported}).

\subsection{Federated Learning Path}

The FL pathway is specified in terms of communication rounds $r=1,\ldots,R$,
local epoch count $\tau$, and the FedAvg update rule (Eq.~\ref{eq:fedavg}).
Privacy amplification via local differential privacy is applied by clipping
client gradients to $\ell_2$ norm $C$ and adding calibrated Gaussian noise
$\mathcal{N}(\mathbf{0},\; \sigma_{\mathrm{dp}}^2 C^2 \mathbf{I})$
before aggregation, with the privacy budget characterized by
$(\varepsilon,\delta)$-DP~\cite{kairouz2021advances}.
The noise level satisfies
\begin{equation}
  \sigma_{\mathrm{dp}} \geq
  \frac{C\sqrt{2\ln(1.25/\delta)}}{\varepsilon \sqrt{n_k}},
  \label{eq:dp_noise}
\end{equation}
where $n_k$ is the number of local training samples used per round.
Flower or an equivalent FL orchestration framework is planned for
deployment; no federated training or aggregation experiments are reported
in this paper.
